\begin{trapbox}
	\begin{description}[style=nextline, leftmargin=2.8em, labelsep=0.5em, itemsep=0.35em]
		\item[\textbf{\textcolor{CTrap}{易错点一（非零条件）}}] 使用结论时，疏忽$f(x)\neq0$的前提，直接套周期公式。
		\item[\textbf{\textcolor{CTrap}{正确理解（非零要求）：}}] 必须满足$f(x)\neq0$，否则倒数无意义，周期结论不成立。
		\item[\textbf{\textcolor{CTrap}{错因分析（分母为零）：}}] 倒数运算要求分母非零，若$f(x)=0$，条件$f(x+a)=1/f(x)$无定义，结论不适用。
	\end{description}

	\medskip
	\begin{description}[style=nextline, leftmargin=2.8em, labelsep=0.5em, itemsep=0.35em]
		\item[\textbf{\textcolor{CTrap}{易错点二（周期误判）}}] 只看到一次变换，认为周期是$a$，错误地把$f(x+a)=1/f(x)$当作周期关系。
		\item[\textbf{\textcolor{CTrap}{正确理解（周期二倍）：}}] 必须迭代两次：$f(x+2a)=\frac{1}{f(x+a)}$，代入条件得 $f(x+2a)=f(x)$，因此周期是$2a$。
		\item[\textbf{\textcolor{CTrap}{错因分析（缺少迭代）：}}] 仅观察$f(x+a)$与$f(x)$的倒数关系，没有继续推导$x+2a$处的情形。
	\end{description}

	\medskip
	\begin{description}[style=nextline, leftmargin=2.8em, labelsep=0.5em, itemsep=0.35em]
		\item[\textbf{\textcolor{CTrap}{易错点三（负号遗漏）}}] 在$f(x+a)=-\frac{1}{f(x)}$时，推导中忘记两个负号抵消，得出周期$T=4a$等错误结论。
		\item[\textbf{\textcolor{CTrap}{正确理解（负号抵消）：}}] 类似地，$f(x+2a)=-\frac{1}{f(x+a)}$，代入 $f(x+a)=-\frac{1}{f(x)}$ 得 $f(x+2a)=f(x)$，周期仍为$2a$。
		\item[\textbf{\textcolor{CTrap}{错因分析（符号粗心）：}}] 对复合分式计算粗心，未正确处理负号的叠乘。
	\end{description}
\end{trapbox}
